\section{Energy and a limiting configuration}\label{sec:energy}

The sampling statements in Proposition~\ref{prop:sampling} concern finite
tables alone.  In this section they will be applied to functions obtained
from an isometric action chosen after those tables are known.  The order
is important: we first obtain a bound on every generator displacement
using the pointwise relative counts in the six-letter table.  Only then
do we use total variation estimates for action-dependent functions.

We use the following standard notation for a complete $\CAT(0)$ space
$X$.  Its completed tangent cone at $x$ is denoted by $T_xX$, with vertex
$0$.  For $y\in X$, the vector $\log_x y$ has length $d(x,y)$ and the
direction of the geodesic from $x$ to $y$.  In a Euclidean metric cone set
\[
 \ip{u}{v}=\frac{\norm{u}^2+\norm{v}^2-d(u,v)^2}{2}.
\]
Thus $\ip{u}{v}=\norm{u}\norm{v}\cos\angle(u,v)$, where the angle is
truncated at $\pi$.  Multiplication of a vector by a nonnegative real
number always means radial dilation.  No additive vector-space structure
on a cone is assumed.

We will use strong convexity of squared distance, first variation, the
flat-triangle equality theorem, and the tangent-cone theorem for
$\CAT(0)$ spaces; see \cite[Chapters~II.1--II.3, in particular
II.2.9--II.2.10, II.3.6, and II.3.19]{BH1999}.  These results do not require
properness.  In particular, comparison of angles with comparison angles
gives
\begin{equation}\label{energy:log-comparison}
 d(y,z)^2\ge
 d(x,y)^2+d(x,z)^2
 -2\ip{\log_x y}{\log_x z}.
\end{equation}
All the tangent cones below are completed cones.

\subsection{A minimizing basepoint and uniform displacement bounds}

The first-variation characterization of an energy minimizer in a
tangent cone is part of the harmonic-map framework of
Gromov~\cite[Sections~3.3 and~3.7]{Gromov2003} and
Izeki--Nayatani~\cite{IzekiNayatani2005}.
We give the argument together with the existence of a minimizing
basepoint under the action hypotheses used here.

\begin{lemma}\label{energy:minimizer}
Let an infinite finitely generated group $G$ act properly and
cocompactly by isometries on a proper complete $\CAT(0)$ space $X$.
For a finite symmetric generating alphabet $S$, counted with its given
multiplicities, the function
\[
 E(x)=\frac1{|S|}\sum_{s\in S}d(x,sx)^2
\]
has a minimizer and its minimum is positive.  At any minimizing
basepoint, the vectors $v_s=\log_x(sx)$ satisfy
\begin{equation}\label{energy:zero-barycenter}
 \frac1{|S|}\sum_{s\in S}\ip{h}{v_s}\le0
 \quad(h\in T_xX).
\end{equation}
Equivalently, $0$ is their barycenter, and
\begin{equation}\label{energy:variance}
 \frac1{|S|}\sum_{s\in S}d(h,v_s)^2
 \ge \norm{h}^2+\frac1{|S|}\sum_{s\in S}\norm{v_s}^2.
\end{equation}
\end{lemma}

\begin{proof}
Choose a minimizing sequence $x_j$ and a compact set $C$ with $GC=X$.
There are $g_j\in G$ such that $y_j=g_jx_j\in C$.  Bounded energy bounds
each individual displacement $d(x_j,sx_j)$, since $S$ is finite.  It
therefore bounds
\[
 d(y_j,(g_jsg_j^{-1})y_j).
\]
The closed neighborhood of $C$ of any fixed finite radius is compact,
by properness of $X$.  Properness of the action implies that only
finitely many group elements can carry a point of $C$ into that
neighborhood.  For each $s$ there are consequently only finitely many
possibilities for $g_jsg_j^{-1}$.  Passing to a subsequence, the entire
conjugated generating tuple is constant and $y_j$ converges to a point
$y\in C$.  Its energy for that tuple is the original infimum.  Fixing
one element $g_{j_0}$ in the subsequence, the point $g_{j_0}^{-1}y$
attains the infimum for the original tuple.  An energy of zero would
make all generators, and hence $G$, fix a point.  A proper action has
finite point stabilizers, so the minimum is positive.

Fix a minimizing point $x$.  First let $h$ be a tangent vector
represented by a geodesic $x_t$ issuing from $x$, parameterized with
initial speed $\norm h$.  For sufficiently small $t>0$, convexity at
the midpoints of the crossed geodesics gives
\[
 d(x_t,sx_t)^2\le
 \frac12d(x_{2t},sx)^2+\frac12d(x,sx_{2t})^2.
\]
After averaging over the symmetric alphabet, the two terms on the
right have the same average: apply $s^{-1}$ to the second distance and
replace $s$ by $s^{-1}$.  Minimality and first variation now imply
\[
 0\le
 \left.\frac{d}{dt}\right|_{0+}
 \frac1{|S|}\sum_s d(x_{2t},sx)^2
 =-\frac4{|S|}\sum_s\ip{h}{v_s}.
\]
Represented vectors are dense in the completed tangent cone, so
continuity gives Equation~\eqref{energy:zero-barycenter} for every $h$.
The cone cosine formula immediately gives
Equation~\eqref{energy:variance}; in particular $0$ minimizes the mean
squared distance to the $v_s$.
\end{proof}

For each aspherical instance $P$ of Theorem~\ref{thm:diagram}, its group
$G$ is infinite.  Indeed, the parity quotient in Lemma~\ref{lem:model-presentation}
makes $G$ nontrivial.  If $G$ were finite, the universal cover would be a
finite contractible complex and
\[
 1=\chi(\widetilde P)=|G|\chi(P),
\]
which forces $|G|=1$.  Thus Lemma~\ref{energy:minimizer} applies to any
geometric action considered here.  Rescale its metric so that at the
minimizing basepoint
\begin{equation}\label{energy:normalization}
 \frac1N\sum_{s\in S_n}\norm{v_s}^2=1.
\end{equation}
Rescaling preserves completeness, properness, and the $\CAT(0)$
condition.

\begin{lemma}\label{energy:displacement-bound}
Suppose the six-letter table has pointwise relative error at most
$\varepsilon<1$ in the counts for positions $1$ and $k$, for every
$2\le k\le6$.  Then every action normalized as in
Equation~\eqref{energy:normalization} satisfies
\begin{equation}\label{energy:max-displacement}
 \max_{s\in S_n}\norm{v_s}^2
 \le25\frac{1+\varepsilon}{1-\varepsilon}.
\end{equation}
The bound is independent of the action.
\end{lemma}

\begin{proof}
Write $Q_{1k}(a,b)$ for the number of included rows whose first
coordinate is $a$ and whose $k$th coordinate is $b$.  The asserted
relative estimate is
\[
 (1-\varepsilon)N^{1-\delta}
 \le Q_{1k}(a,b)\le
 (1+\varepsilon)N^{1-\delta}.
\]
Summing over $b$ shows that rows beginning with any fixed $a$ are
nonempty and that their conditional $k$th-coordinate distribution is
bounded above, pointwise, by
\[
 \frac{1+\varepsilon}{1-\varepsilon}\frac1N.
\]
The relation attached to a row $(a,s_2,\ldots,s_6)$ gives
$a^{-1}=s_2\cdots s_6$.  Since
$d(x,a^{-1}x)=d(x,ax)$, the triangle inequality and Cauchy--Schwarz give
\[
 \norm{v_a}^2\le5\sum_{k=2}^6\norm{v_{s_k}}^2.
\]
Average over the rows beginning with $a$ and use the preceding
pointwise bound and Equation~\eqref{energy:normalization}.  Each of the
five conditional averages is at most
$(1+\varepsilon)/(1-\varepsilon)$, proving the assertion.  This uses
only nonnegative weights, and therefore holds even when the weights
were chosen from an action after the table was revealed.
\end{proof}

\subsection{Conditional tests in tangent cones}

For a geodesic from $u$ to $v$ in a $\CAT(0)$ Euclidean cone, write
$[u,v]_t$ for its point at parameter $t$.  Its squared
radial norm satisfies the exact identity
\begin{equation}\label{energy:radial-identity}
 \norm{[u,v]_t}^2
 =(1-t)\norm u^2+t\norm v^2-t(1-t)d(u,v)^2.
\end{equation}
If either endpoint is the cone vertex, the identity is immediate.
Otherwise the segment is a Euclidean segment in the sector between
its two rays when their angular distance is less than $\pi$.  When
their angular distance is at least $\pi$, it passes through the cone vertex, and
$\norm{[u,v]_t}=|(1-t)\norm u-t\norm v|$ gives the same identity.
These two cases follow directly from the Euclidean-cone metric;
see \cite[Definition~I.5.6 and Proposition~I.5.10(1)--(2)]{BH1999}.

The scalar-product concavity used next is recorded by
Gigli--Nobili~\cite[Proposition~2.6, equation~(2.7f), and the proof of
Proposition~3.8]{GigliNobili2021}. The inductive-mean estimate is the
unit-variance case of Sturm's bound~\cite[Definition~4.6 and
Theorem~4.7, proof~\textup{(a)}]{Sturm2003}. We include the short proofs in the present
normalization.

\begin{lemma}\label{energy:cone-mean}
In a complete $\CAT(0)$ metric cone, $u\mapsto\ip{h}{u}$ is concave
along every geodesic, for fixed $h$.  Suppose a finite probability
distribution of vectors $V$ has barycenter $0$ and
$\mathbb E\norm V^2=1$.  If $V_1,\ldots,V_t$ are independent samples
and
\[
 m_1=V_1,\qquad
 m_j=[m_{j-1},V_j]_{1/j},
\]
then
\begin{equation}\label{energy:inductive-mean}
 \mathbb E\norm{m_t}^2\le\frac1t.
\end{equation}
For any deterministic list $V_1,\ldots,V_t$ one also has
\begin{equation}\label{energy:mean-concavity}
 \frac1t\sum_{j=1}^t\ip{h}{V_j}\le\ip{h}{m_t}.
\end{equation}
\end{lemma}

\begin{proof}
Subtract strong convexity of $d(h,\cdot)^2$ from
Equation~\eqref{energy:radial-identity} in the identity
$2\ip{h}{u}=\norm h^2+\norm u^2-d(h,u)^2$.  The quadratic terms
cancel, proving concavity.  Induction gives
Equation~\eqref{energy:mean-concavity}.

For the independent samples, condition on $m_{j-1}=u$.
The barycenter variance inequality gives
$\mathbb E[d(u,V_j)^2\mid u]\ge\norm u^2+1$.  Strong convexity
of squared distance from the cone vertex therefore gives, with
$a_j=\mathbb E\norm{m_j}^2$,
\[
 a_j\le(1-1/j)^2a_{j-1}+1/j^2.
\]
Starting with $a_1=1$, induction yields $a_j\le1/j$.
\end{proof}

Consider now a sequence of table instances indexed by $\nu$ for which
the conclusions of Proposition~\ref{prop:sampling} hold with errors
tending to zero, for every fixed test and every fixed positive integer
$t$.  Suppose their groups have geometric actions as above, and choose
the normalized minimizing basepoints $x_\nu$.  Write
$\Omega_\nu=\Omega_{n_\nu}$ for the nonempty finite set of master
matches and $\mu_\nu=\mu_{N_\nu}$ for its uniform distribution.
Lemma~\ref{energy:displacement-bound} gives an
absolute constant $M$ such that
\begin{equation}\label{energy:uniform-bound}
 \norm{v_s}\le M
 \quad\text{for all letters in all sufficiently large instances.}
\end{equation}
All statements about the limit are unchanged by discarding finitely
many instances.  The actions, their spaces, and the functions in the
next lemma may depend on the complete table instances.

\newpage
\begin{lemma}[Uniform conditional test]\label{lem:conditional-test}
Let $I$ and $\{e\}$ be disjoint sets of fields with admissible union.
For each $\nu$, let $h_I$ be a function of the labels on $I$ with values
in $T_{x_\nu}X_\nu$, and suppose $\norm{h_I}\le B$ for a constant
independent of $\nu$ and the match.  For either fixed choice
$\sigma(s)=s$ or $\sigma(s)=s^{-1}$,
\begin{equation}\label{energy:conditional-test}
 \limsup_{\nu\to\infty}
 \mathbb E_{\mu_\nu}
 \ip{h_I}{v_{\sigma(e)}}\le0.
\end{equation}
Here $e$ in a subscript denotes the letter occupying that field.
\end{lemma}

\begin{proof}
Fix $t$.  Use the coupling in Proposition~\ref{prop:sampling}: it
gives $t$ matches agreeing on $I$, each branch marginal at total
variation distance at most $\eta_{\nu,t}$ from $\mu_\nu$, and their
$e$-labels at distance at most $\eta_{\nu,t}$ from independent uniform
letters, where $\eta_{\nu,t}\to0$.  Increasing the error if necessary
allows one bound for all these comparisons.  Applying the fixed
inversion $\sigma$ preserves this assertion.

The vector $h_I$ has the same value in every branch.  Form the
inductive mean $m_t$ of the branch vectors $v_{\sigma(e_j)}$.
Equation~\eqref{energy:mean-concavity} gives, for every coupled tuple,
\[
 \frac1t\sum_{j=1}^t\ip{h_I}{v_{\sigma(e_j)}}
 \le\ip{h_I}{m_t}\le B\norm{m_t}.
\]
The last step uses only the norm bound; it does not assume that $h_I$
is independent of the leaves.

Under independent uniform leaves,
Equation~\eqref{energy:inductive-mean} applies by
Equations~\eqref{energy:variance} and~\eqref{energy:normalization}.
Under the coupling, $\norm{m_t}\le M$, by convexity of the ball of
radius $M$.  Consequently the total variation comparison changes the
expected squared norm by at most $2M^2\eta_{\nu,t}$.  The branch
comparisons change the expectation of each scalar-product test by at
most $2BM\eta_{\nu,t}$.  It follows that
\[
 \mathbb E_{\mu_\nu}\ip{h_I}{v_{\sigma(e)}}
 \le B\sqrt{1/t+2M^2\eta_{\nu,t}}+2BM\eta_{\nu,t}.
\]
First take the upper limit in $\nu$ with $t$ fixed, and then let
$t\to\infty$.  All error bounds used only $B$ and $M$, so the
conclusion is uniform over the action-dependent functions in the
statement.
\end{proof}

\subsection{Products of orbit configurations}

Enlarge $\Gamma$ by adding, for each of its prescribed two-edge paths,
the complementary two-edge path of its decoration square.  Denote
this fixed finite graph by $\widehat\Gamma$; its additional midpoint
vertices are distinct abstract vertices.  A master match gives a
consistent orbit realization of this graph in $X_\nu$: choose one
anchor at $x_\nu$, and send a vertex reached by the successive letters
$s_1,\ldots,s_k$ to $s_1\cdots s_kx_\nu$.  All core closed walks are
relations, and the decoration square relations guarantee consistency
on the added paths.  Fixing one path for each vertex shows that these
orbit points are at distance at most $LM$ from $x_\nu$, where $L$
depends only on $\widehat\Gamma$.

Take the finite product over the matches with metric
\begin{equation}\label{energy:product-metric}
 d_\nu((y_m),(z_m))^2
 =\frac1{|\Omega_\nu|}\sum_{m\in\Omega_\nu}d(y_m,z_m)^2,
\end{equation}
based at the tuple of minimizing basepoints.  It is a complete
$\CAT(0)$ space.  The vertex tuples have bounded distance from this
basepoint.  Fix a nonprincipal ultrafilter on the sequence index set
$\mathbb N$ and take the pointed metric ultralimit $Z$.  This space is
complete by \cite[Lemma~I.5.53]{BH1999}, and it is $\CAT(0)$ by
\cite[Corollary~II.3.10(2)]{BH1999}.  The bounded vertex tuples give a
configuration in $Z$ indexed by the vertices of $\widehat\Gamma$.
Midpoints
and fixed-fraction points on segments in the products are formed
coordinatewise, and their ultralimits are the corresponding points
in $Z$ by uniqueness of geodesics.  We identify the abstract vertex
names with the resulting points.  Neither properness nor a dimension
bound for $Z$ is needed.

\newpage
\begin{lemma}\label{energy:path-bounds}
Every edge in $\widehat\Gamma$ has length one in $Z$.  Every two-step
path using two distinct fields of one decoration square has squared
endpoint distance at least two.  More generally, a path of $k\le3$
distinct core fields whose field set is admissible has squared
endpoint distance at least $k$.
\end{lemma}

\begin{proof}
For an edge with letter $s$, translation at its initial orbit vertex
identifies its squared length with $\norm{v_s}^2$, or with the same
quantity for $s^{-1}$.  Single-field total variation convergence,
Equation~\eqref{energy:uniform-bound}, and the normalization show that
its squared product length tends to one.

For a path with oriented readings $a_1,\ldots,a_k$, translate its
penultimate orbit vertex to $x_\nu$.  Its earlier endpoint then becomes
$a_{k-1}^{-1}\cdots a_1^{-1}x_\nu$, and its next endpoint becomes
$a_kx_\nu$.  Set
\[
 h_I=\log_{x_\nu}(a_{k-1}^{-1}\cdots a_1^{-1}x_\nu).
\]
This is a function of precisely the preceding labels, and
$\norm{h_I}\le(k-1)M$.  Equation~\eqref{energy:log-comparison} gives
\[
 d(u_0,u_k)^2\ge d(u_0,u_{k-1})^2+\norm{v_{a_k}}^2
       -2\ip{h_I}{v_{a_k}}
\]
in every factor.  Average, apply Lemma~\ref{lem:conditional-test},
and take the ultralimit.  Induction adds at least one at each step.
Every prefix of the field sets in the statement is admissible, so
the conditional test is available at each step.  Distinct fields
here need not carry distinct numerical letter values.
\end{proof}

\subsection{Quadrilateral equality and hexagon centers}

The path estimates give lower bounds on squared distances.  We now
compare them with a $\CAT(0)$ upper bound on the sum of the two
diagonal squares of a quadrilateral.  Equality will first determine
the diagonals of every decoration square, and then force the three
opposite pairs in each core hexagon to share a midpoint.

\begin{lemma}\label{energy:quadrilateral}
For four points $z_1,z_2,z_3,z_4$ in a $\CAT(0)$ space, let
$m=\operatorname{mid}(z_1,z_3)$ and
$n=\operatorname{mid}(z_2,z_4)$.  Then
\begin{equation}\label{energy:quadrilateral-inequality}
 d(z_1,z_3)^2+d(z_2,z_4)^2+4d(m,n)^2
 \le\sum_{i=1}^4d(z_i,z_{i+1})^2,
 \qquad z_5=z_1.
\end{equation}
In particular, equality between the two diagonal-square sum and the
side-square sum forces the diagonal midpoints to coincide.
\end{lemma}

\begin{proof}
Apply midpoint convexity to $d(m,z_2)^2$ and $d(m,z_4)^2$ using the
segment $z_1z_3$, and then to $d(m,n)^2$ using $z_2z_4$.  Substituting
the first two inequalities into the third and multiplying by four
gives Equation~\eqref{energy:quadrilateral-inequality}.
\end{proof}

Each decoration square has unit sides.  Both of its diagonals have
squared length at least two by Lemma~\ref{energy:path-bounds}.
Lemma~\ref{energy:quadrilateral} forces both squared lengths to equal
two.  In particular,
\begin{equation}\label{energy:two-step-equality}
 d(u,w)^2=2
 \quad\text{for every nonbacktracking two-edge core path }u,v,w.
\end{equation}
Every such path was assigned a decoration table, including paths
whose two fields lie in the same fan.

Recall that a core hexagon is obtained from three
noncollinear Fano points $i,j,k$, their joining lines, and a choice of
$X$ or an available $Y$ representative at each of the three points,
with all six required edges present.  The squared distances of its
vertices $z_0,\ldots,z_5$ at cyclic separation one and two are,
respectively, one and two.  Each three-edge half-path uses at most
two fields of any one fan, so its opposite endpoints have squared
distance at least three.

Apply Lemma~\ref{energy:quadrilateral} to
$(z_0,z_1,z_3,z_4)$.  The side-square sum is $1+2+1+2=6$, whereas each
diagonal-square is at least three.  Both are therefore three and
their midpoints coincide.  Applying the same argument to
$(z_1,z_2,z_4,z_5)$ identifies the third opposite-pair midpoint with
the first two.  Denote the resulting center of this hexagon by $c$.
Every vertex has distance
\begin{equation}\label{energy:radius}
 r=\frac{\sqrt3}{2}
\end{equation}
from $c$.

\begin{lemma}\label{energy:hexagon-flatness}
For a core hexagon, every triangle formed by its center and two of
its vertices is Euclidean flat, including the degenerate case of
opposite vertices.  The corresponding tangent vectors satisfy
\begin{equation}\label{energy:hexagon-products}
 \ip{\log_c z_i}{\log_c z_j}
 =r^2-\frac12d(z_i,z_j)^2.
\end{equation}
\end{lemma}

\begin{proof}
The opposite pairs are geodesic segments through $c$, so their
directions in $T_cZ$ are antipodal and their radial lengths are $r$.
If $v^+$ and $v^-$ are any such opposite vectors, the angular
triangle inequality implies
\begin{equation}\label{energy:antipodal-products}
 \ip{u}{v^+}+\ip{u}{v^-}\le0
 \quad(u\in T_cZ).
\end{equation}
Indeed, the two angles from $u$ to their directions have sum at
least $\pi$, and the sum of their cosines is nonpositive.

Set $u=\log_c z_i$, $v^+=\log_c z_j$, and
$v^-=\log_c z_{j+3}$, with indices modulo six.  Log comparison gives
\[
 \ip{u}{v^+}\ge r^2-\tfrac12d(z_i,z_j)^2,
 \qquad
 \ip{u}{v^-}\ge r^2-\tfrac12d(z_i,z_{j+3})^2.
\]
The two prescribed squared distances sum to $3=4r^2$, so the lower
bounds sum to zero.  Equation~\eqref{energy:antipodal-products}
forces equality in each bound.  Thus the angle at $c$ equals its
Euclidean comparison angle.  For a nondegenerate triangle the
flat-triangle equality theorem gives a flat triangle; for an
opposite or repeated pair the statement is a segment or a point.
\end{proof}

\begin{lemma}\label{energy:common-center}
All core hexagons have the same center $c$.  Every core vertex lies
in a core hexagon, and hence is at distance $r$ from $c$.
\end{lemma}

\begin{proof}
First use only $X$ representatives.  A noncollinear ordered triple
of Fano points is a basis of $\mathbb F_2^3$.  Reordering a basis
does not change its hexagon, and replacing one basis vector $i$ by
$i+j$ preserves the line $ij$ and the third point $k$.  The old and
new hexagons therefore share the opposite pair $(X_k,ij)$, and their
centers are its common midpoint.  Permutations and elementary
additions generate all changes of basis over $\mathbb F_2$, so the
centers of all $X$ hexagons agree.

Replace one $Y$ representative in a core hexagon by its $X$
representative.  The new hexagon is a core hexagon, since every corresponding
$X$ edge exists.  Either of the unchanged point representatives,
together with its opposite line, is an opposite pair shared by the
two hexagons.  Their centers agree.  Repeating this replacement
connects every core hexagon to one using only $X$ representatives.

Every $X_i$ lies in a hexagon obtained by extending $i$ to a basis.
Every line lies in one by choosing two points spanning that line
and a third point outside it.  Each $Y_i$ has two distinct incident
lines available.  Choose a point other than $i$ on each of them;
these two points together with $i$ are noncollinear and yield a
core hexagon containing $Y_i$.  This includes $Y_D$, using its two
full lines.  The radius assertion follows.
\end{proof}

\subsection{The ternary midpoint test}

The square and hexagon comparisons have supplied a common center.
We now use three admissible core fields to control one further angle: the angle
at $CD$ between the segment to the midpoint of $X_C,Y_C$ and the loose
edge to $X_D$.

Put
\[
 H_0=CD,\qquad p=X_C,\qquad q=Y_C,\qquad z=X_D,
 \qquad J=\operatorname{mid}(p,q)
\]
in $Z$.  The following argument gives an angle statement in $Z$
without any assumption that angles commute with ultralimits.

Choose another line $H\ne H_0$ through $C$.  The core four-cycle
$(p,H_0,q,H)$ has unit sides and both diagonals of squared length
two, by Equation~\eqref{energy:two-step-equality}.  Its diagonal
midpoints coincide by Lemma~\ref{energy:quadrilateral}, so
\begin{equation}\label{energy:ternary-radii}
 J=\operatorname{mid}(H_0,H),\qquad
 d(H_0,J)^2=\frac12,\qquad d(H_0,z)=1.
\end{equation}

In each orbit factor let $J_m$ be the midpoint of the realizations
of $p$ and $q$.  Translate the realization of $H_0$ to $x_\nu$.
The realizations of $p,q$ are then given by the two letters on the
adjacent $F_C$ edges, with their fixed orientation signs.  Therefore
the translated vector toward $J_m$ is a function $h_I$ of just
those two fields, with norm at most $M$.  The translated vector
toward $z$ is $v_{\sigma(e_0)}$, where $e_0$ is the loose field.
These three fields are admissible.  Lemma~\ref{lem:conditional-test}
shows that the upper limit of the averaged scalar product of these
two vectors is at most zero.

Fix $a,b\in(0,1]$.  In each factor take the points at fractions
$a,b$ along the segments from $H_0$ to $J_m,z$.  Their tangent
vectors are the corresponding scalar multiples of the two vectors
just considered.  Apply Equation~\eqref{energy:log-comparison},
average over factors, and take the ultralimit.  Denoting the
limiting fractional points by $J_a,z_b$, respectively, gives
\begin{equation}\label{energy:fractional-angle}
 d(J_a,z_b)^2\ge
 a^2d(H_0,J)^2+b^2d(H_0,z)^2
 =\frac{a^2}{2}+b^2.
\end{equation}
Thus every such comparison angle at $H_0$ is at least $\pi/2$.
Letting both fractions decrease to zero in $Z$ proves
\begin{equation}\label{energy:ternary-angle}
 \angle_{H_0}(J,z)\ge\frac\pi2.
\end{equation}
Only distances of fixed-fraction points passed through the
ultralimit; the angle was computed afterwards in the limiting
space.

\subsection{The resulting finite configuration}

We collect the conclusions in the form needed in
Section~\ref{sec:obstruction}.

\begin{proposition}\label{prop:configuration}
Let $P_\nu$ be a sequence of aspherical instances whose master-match
statistics and six-letter relative counts satisfy the conclusions
of Proposition~\ref{prop:sampling}, with errors tending to zero for
every fixed test and every fixed $t$.  If every $\pi_1(P_\nu)$ has
a proper cocompact isometric action on a proper complete $\CAT(0)$
space, then there are a complete $\CAT(0)$ space $Z$, a
configuration indexed by the core vertices of $\Gamma$, and a
point $c\in Z$ with the following properties.
\begin{enumerate}
\item Every core edge has length one.  Every nonbacktracking
two-edge core path has squared endpoint distance two.  Every
three-edge core path with distinct admissible fields has squared
endpoint distance at least three.
\item Every core hexagon has its three opposite-pair midpoints
equal to $c$.  All core vertices have distance $r=\sqrt3/2$ from
$c$.  For any two vertices occurring in one core hexagon, their
central triangle with $c$ is flat and their tangent scalar
product is given by Equation~\eqref{energy:hexagon-products}.
\item With $H_0=CD$ and $J=\operatorname{mid}(X_C,Y_C)$, one has
\[
 d(H_0,J)^2=\frac12,\qquad d(H_0,X_D)=1,
 \qquad \angle_{H_0}(J,X_D)\ge\frac\pi2.
\]
More precisely, all the fractional-point inequalities in
Equation~\eqref{energy:fractional-angle} hold.
\end{enumerate}
\end{proposition}

\begin{proof}
Use the normalized minimizing actions of
Lemma~\ref{energy:minimizer} and the uniform bound of
Lemma~\ref{energy:displacement-bound}.  The product construction,
Lemma~\ref{energy:path-bounds}, and the decoration-square argument
give the first assertion.  Lemmas~\ref{energy:hexagon-flatness}
and~\ref{energy:common-center} give the second.  The fixed-fraction
argument above gives the third.
\end{proof}
